\SourceLeaf{438}%
\[
\begin{aligned}
\varrho d\xi&=Xds\\
\varrho d\eta&=Yds\\
\varrho d\zeta&=Zds
\end{aligned}
\]
und wir haben hier anstatt $4$ krummer Linien auf der Oberfläche der Hülfskugel nur $3$ zu betrachten. Jedes Element der zweiten Linie ist also als in dem grössten Kreise $L\lambda$ liegend anzusehen, und der positive oder negative Werth von $\varrho$ wird sich auf die Concavität oder Convexität der Curve nach der Richtung der Normale beziehen.

\begin{center}14.\end{center}
\noindent\hspace*{1.2em}Wir wollen nunmehro den sphärischen Winkel auf der Hülfskugel untersuchen, den der grösste Kreis von $L$ nach $\lambda$ mit demjenigen macht, der von $L$ nach einem der festen Punkte $(1),(2),(3)$, z.\,B. nach $(3)$, geht. Um hierin etwas bestimmtes zu haben, zählen wir von $L(3)$ nach $L\lambda$ in demjenigen Sinn herum, in welchem $(1),(2)$ und $(3)$ liegen. Nennen wir $\varphi$ diesen Winkel, so ist in Folge des Satzes Art. $7$
\[
\sin L(3).\sin L\lambda.\sin\varphi=Y\xi-X\eta
\]
oder, da $L\lambda=90^{0}$ und
\[
\sin L(3)=\BareRoot{(XX+YY)}=\BareRoot{(1-ZZ)}:
\]
\[
\sin\varphi=\frac{Y\xi-X\eta}{\BareRoot{(XX+YY)}}.
\]
Es ist ferner
\[
\sin L(3).\sin L\lambda.\cos\varphi=\zeta
\]
oder
\[
\cos\varphi=\frac{\zeta}{\BareRoot{(XX+YY)}}
\]
und
\[
\operatorname{tang}\varphi=\frac{Y\xi-X\eta}{\zeta}=\frac{\zeta'}{\zeta}.
\]
Man hat also
\[
d\varphi=\frac{\zeta Yd\xi-\zeta Xd\eta-(Y\xi-X\eta)d\zeta+\xi\zeta dY-\eta\zeta dX}{(Y\xi-X\eta)^{2}+\zeta\zeta}.
\]
